\section*{Chapitre \ref{ch:g11:antibiotic-resistance} --- Les bactéries et la résistance aux antibiotiques}

\begin{solution}{exo:g11:antibiotic-resistance:1}
Les \omterm{def:g11:mutations:mutation}{mutations} qui surviennent au cours des divisions du clone, et le
\omterm{prop:g11:antibiotic-resistance:variation}{transfert horizontal de gènes} --- de l'\omterm{def:g10:universal-dna:information}{ADN}, le plus souvent un
plasmide, reçu d'une autre bactérie.
\end{solution}

\begin{solution}{exo:g11:antibiotic-resistance:2}
Une molécule qui tue les bactéries ou arrête leur croissance en
agissant sur une structure que les bactéries possèdent et que les
\omterm{def:g10:cells-common-unit:cell}{cellules} humaines n'ont pas (la paroi, le \omterm{def:g10:cells-common-unit:organelle}{ribosome} bactérien, les
\omterm{def:g11:enzymes-and-phenotype:enzyme}{enzymes} bactériennes de réplication). Un rhume est viral, et les virus
n'ont aucune de ces cibles.
\end{solution}

\begin{solution}{exo:g11:antibiotic-resistance:3}
Une \omterm{def:g11:enzymes-and-phenotype:enzyme}{enzyme} qui détruit l'\omterm{def:g11:antibiotic-resistance:antibiotic}{antibiotique}~; une cible mutée sur laquelle le
médicament ne se fixe plus~; une pompe qui expulse le médicament (ou
une paroi qui l'exclut).
\end{solution}

\begin{solution}{exo:g11:antibiotic-resistance:4}
Sensible (20 au-dessus de 18), mais de justesse. D, dont la marge
au-dessus du seuil est la plus large, puis A.
\end{solution}

\begin{solution}{exo:g11:antibiotic-resistance:5}
Non. Sur les répliques sur velours, les colonies résistantes
apparaissent aux mêmes positions sur chaque boîte à \omterm{def:g11:antibiotic-resistance:antibiotic}{antibiotique}~: les
\omterm{def:g10:cells-common-unit:cell}{cellules} résistantes existaient donc dans les colonies d'origine, qui
n'avaient jamais rencontré le médicament. L'\omterm{def:g11:antibiotic-resistance:antibiotic}{antibiotique} révèle et
sélectionne~; il n'induit pas.
\end{solution}

\begin{solution}{exo:g11:antibiotic-resistance:6}
30 divisions~: $2^{30} \approx 10^9$ \omterm{def:g10:cells-common-unit:cell}{cellules}. Environ $10^9$ divisions
les ont produites~: une dizaine de \omterm{def:g10:cells-common-unit:cell}{cellules} résistantes en moyenne ---
davantage si la \omterm{def:g11:mutations:mutation}{mutation} est survenue tôt.
\end{solution}

\begin{solution}{exo:g11:antibiotic-resistance:7}
Vers le jour 3, quand les \omterm{def:g10:cells-common-unit:cell}{cellules} sensibles sont tombées à environ
0.2~\% et les résistantes montées à environ 1~\%. Au jour 6, le total
en pointillés est revenu à 100~\%, fait presque uniquement de \omterm{def:g10:cells-common-unit:cell}{cellules}
résistantes (les sensibles survivantes du jour 3 ont repoussé aussi,
mais à partir d'une base plus petite).
\end{solution}

\begin{solution}{exo:g11:antibiotic-resistance:8}
Un \omterm{def:g10:universal-dna:gene}{gène} chromosomique ne passe qu'aux descendantes du clone~; un
plasmide passe aussi horizontalement, à des \omterm{def:g10:cells-common-unit:cell}{cellules} sans lien de
parenté et à d'autres \omterm{def:g10:biodiversity-scales:species}{espèces}, et un seul transfert donne à une \omterm{def:g10:cells-common-unit:cell}{cellule}
la résistance d'un coup, sans attendre une \omterm{def:g11:mutations:mutation}{mutation}.
\end{solution}

\begin{solution}{exo:g11:antibiotic-resistance:9}
Au jour 3, les \omterm{def:g10:cells-common-unit:cell}{cellules} sensibles sont tombées à une fraction de pour
cent tandis que les résistantes montent~; sans le médicament, toutes
les survivantes repoussent, et la population qui revient est bien plus
riche en \omterm{def:g10:cells-common-unit:cell}{cellules} résistantes que l'originale. La rechute ne répond
plus au même médicament.
\end{solution}

\begin{solution}{exo:g11:antibiotic-resistance:10}
Les hôpitaux emploient les \omterm{def:g11:antibiotic-resistance:antibiotic}{antibiotiques} intensivement chez de nombreux
patients rassemblés~: la sélection est constante et les souches
sélectionnées passent d'un patient à l'autre par les mains et les
surfaces.
\end{solution}

\begin{solution}{exo:g11:antibiotic-resistance:11}
La streptomycine agit en se fixant sur une \omterm{def:g11:gene-expression:protein}{protéine} ribosomique~; une
\omterm{def:g11:gene-expression:protein}{protéine} changée ne la fixe plus et la \omterm{prop:g11:gene-expression:translation}{traduction} continue. Mais le
changement gêne aussi un peu le travail normal du \omterm{def:g10:cells-common-unit:organelle}{ribosome}, si bien que
la synthèse des \omterm{def:g10:chemistry-of-life:families}{protéines}, et la croissance, sont un peu plus lentes
--- un coût de la résistance.
\end{solution}

\begin{solution}{exo:g11:antibiotic-resistance:12}
$10^{-8} \times 10^{-8} = 10^{-16}$ par division~: avec $10^{12}$
bactéries, il est peu probable qu'une \omterm{def:g10:cells-common-unit:cell}{cellule} porte les deux. Dans la
tuberculose, le traitement dure des mois et les bactéries se comptent
en milliards, si bien que la résistance à un seul médicament est
certaine d'être sélectionnée~; deux médicaments à la fois ne laissent
aucune survivante à sélectionner.
\end{solution}

\begin{solution}{exo:g11:antibiotic-resistance:13}
De faibles doses sélectionnent des souches résistantes dans l'intestin
des animaux sans tuer la population~; les employés les acquièrent par
contact~; des plasmides portent les \omterm{def:g10:universal-dna:gene}{gènes} jusque dans des pathogènes
humains et dans l'environnement par le lisier. La résistance au
médicament, et à des médicaments humains apparentés, monte en médecine.
\end{solution}

\begin{solution}{exo:g11:antibiotic-resistance:14}
La résistance coûte souvent quelque chose à la bactérie (un \omterm{def:g10:cells-common-unit:organelle}{ribosome}
plus lent, une \omterm{def:g11:enzymes-and-phenotype:enzyme}{enzyme} à fabriquer)~: sans le médicament, les \omterm{def:g10:cells-common-unit:cell}{cellules}
sensibles croissent un peu plus vite et l'emportent sur les résistantes
au fil des générations. La proportion baisse --- mais les \omterm{def:g10:universal-dna:gene}{gènes} de
résistance restent dans la population à basse fréquence, prêts à être
sélectionnés de nouveau.
\end{solution}

\begin{solution}{exo:g11:antibiotic-resistance:15}
Les bactéries ne s'adaptent pas individuellement. Dans toute population
nombreuse, quelques \omterm{def:g10:cells-common-unit:cell}{cellules} portent, par \omterm{def:g11:mutations:mutation}{mutation} aléatoire ou par un
plasmide reçu d'une autre \omterm{def:g10:cells-common-unit:cell}{cellule}, un \omterm{def:g10:universal-dna:gene}{allèle} qui les rend résistantes.
L'\omterm{def:g11:antibiotic-resistance:antibiotic}{antibiotique} tue les autres~; les \omterm{def:g10:cells-common-unit:cell}{cellules} résistantes se multiplient
et transmettent l'\omterm{def:g10:universal-dna:gene}{allèle} à leurs descendantes. La population devient
résistante parce qu'elle a été triée, non parce que ses membres ont
appris quoi que ce soit.
\end{solution}

\begin{solution}{pb:g11:antibiotic-resistance:1}
\textbf{1.} Environ $10^{10}$ divisions (une par \omterm{def:g10:cells-common-unit:cell}{cellule} produite).

\textbf{2.} $10^{10} \times 10^{-8} = 100$ \omterm{def:g10:cells-common-unit:cell}{cellules} résistantes, en moyenne.

\textbf{3.} Cent \omterm{def:g10:cells-common-unit:cell}{cellules} sur dix milliards sont invisibles sur une
boîte~: le tapis se comporte comme sensible parce que 99.999999~\% de
lui l'est.

\textbf{4.} Amoxicilline~: sensible (24 au-dessus de 17).
Érythromycine~: résistante (12 en dessous de 20). Pénicilline~:
pleinement résistante, aucune inhibition.

\textbf{5.} La pénicilline est sans effet sur cette souche. Le choix
parmi les médicaments efficaces pèse aussi les effets secondaires et le
spectre~; on préfère le médicament efficace au spectre le plus étroit,
pour épargner les autres bactéries du patient.

\textbf{6.} Quatre périodes de 6 heures par jour~: $10^{10} \times 0.1^4 =
10^6$ après 1 jour~; $10^{10} \times 0.1^{12} = 0.01$ après 3 jours ---
autant dire aucune.

\textbf{7.} $100 \times 2^{48} \approx \num{2.8e16}$~: absurde, bien
plus que ce que la plaie pourrait contenir. La croissance est limitée
par la place, les nutriments et le système immunitaire.

\textbf{8.} Passer de 100 \omterm{def:g10:cells-common-unit:cell}{cellules} à $10^{10}$ demande
$\log_2 10^8 \approx 27$ doublements, soit environ 13 heures~; en
pratique la place se libère au cours du premier jour, si bien que du
jour 1 au jour 2 le clone résistant a remplacé la population sensible.

\textbf{9.} Sensibles~: en baisse d'un facteur 10 toutes les 6 heures
jusqu'à disparaître au jour 3. Résistantes~: montant de 100 à remplir
la place en un jour ou deux, puis stables à $10^{10}$. Total~: un creux
pendant le jour 1, puis retour à $10^{10}$, désormais toutes
résistantes.

\textbf{10.} Si le système immunitaire élimine les \omterm{def:g10:cells-common-unit:cell}{cellules} plus vite
que la centaine de résistantes ne peut se multiplier, l'infection
s'achève malgré elles~; si le clone dépasse l'élimination, l'infection
devient résistante. Les \omterm{def:g11:antibiotic-resistance:antibiotic}{antibiotiques} travaillent avec le système
immunitaire, non à sa place.

\textbf{11.} Sensibles après 2 jours~: $10^{10} \times 0.1^8 = 100$.
Résistantes~: environ $10^{10}$ si le clone a rempli la place, ou au
moins plusieurs milliers~: les survivantes sont presque toutes
résistantes.

\textbf{12.} Croître au même rythme conserve la proportion~: la
population est pour ainsi dire résistante à 100~\%.

\textbf{13.} Aucun effet~: le clone résistant ignore l'amoxicilline.

\textbf{14.} La première boîte portait la population d'origine,
sensible à 99.999999~\%~; la boîte de la rechute porte le clone
sélectionné, entièrement résistant, si bien que le tapis pousse
jusqu'au disque.

\textbf{15.} Le \omterm{def:g10:universal-dna:gene}{gène} peut passer à d'autres \omterm{def:g10:biodiversity-scales:species}{espèces} portées par
d'autres patients ou par le personnel, si bien que les autres
pathogènes du service peuvent devenir résistants sans attendre leur
propre \omterm{def:g11:mutations:mutation}{mutation}.

\textbf{16.} Dans le service, 400 cures par an dans une petite
population fermée trient continuellement les bactéries, et les souches
sélectionnées passent d'un patient à l'autre~; en ville, la sélection
est diluée entre de nombreuses personnes et de nombreuses \omterm{def:g10:biodiversity-scales:species}{espèces}.

\textbf{17.} L'antibiogramme~: choisir un médicament auquel la souche
ne résiste pas, de sorte qu'aucun clone ne soit sélectionné. Les cures
complètes~: aucune cure interrompue ne laisse repousser les
survivantes résistantes. L'hygiène~: aucune transmission du clone
sélectionné d'un patient à l'autre. L'isolement~: aucune propagation
horizontale ni verticale à partir des \omterm{def:g11:genes-and-disease:genetic}{porteurs}.

\textbf{18.} Une sélection moindre a laissé les souches sensibles, qui
croissent un peu plus vite, regagner du terrain, et la transmission des
résistantes a été coupée. Elle n'est pas revenue à 5~\% parce que les
souches résistantes et leurs plasmides persistent à basse fréquence et
sont resélectionnés par chaque cure restante.

\textbf{19.} Toute population bactérienne nombreuse contient des
\omterm{def:g10:cells-common-unit:cell}{cellules} résistantes au nouveau médicament par \omterm{def:g11:mutations:mutation}{mutation} fortuite~;
l'utiliser les sélectionne, et en quelques années la souche résistante
est courante, comme pour tous les \omterm{def:g11:antibiotic-resistance:antibiotic}{antibiotiques} introduits jusqu'ici.

\textbf{20.} Environ une centaine de \omterm{def:g10:cells-common-unit:cell}{cellules} résistantes sur dix
milliards avant tout médicament~; le deuxième jour du traitement, elles
ont hérité de la plaie~; arrêter la cure trop tôt a changé l'infection
curable en une infection résistante qui s'est répandue.
\end{solution}
